\section{Real-World Audit: Response Semantics on ImageNet-9}
\label{sec:imagenet9}

Section~\ref{sec:controlled} established the meanings of $E$, $C$, and $F$ in controlled learning problems. We now ask whether the same response semantics describe behavior on natural images, including models that were never trained for our audit. ImageNet-9 is useful for this purpose because its background variants keep the foreground object fixed while changing only the background~\citep{xiao2021noise}. This gives natural factual--counterfactual pairs in which we can ask whether a model's response to background change behaves as evidence, contradiction, or endpoint-null sensitivity.

\textbf{Setup.} We evaluate two complementary groups of models on these same background interventions. The first consists of 24 off-the-shelf ImageNet-1k classifiers. These models were trained independently of our experiment, so they test whether the response semantics established in Section~\ref{sec:controlled} remain meaningful for existing natural-image models. Their 1,000-way predictions are aggregated into the nine ImageNet-9 superclasses before evaluation.

The second group contains 48 models trained directly on ImageNet-9. We fine-tune six backbones---\emph{ResNet-50, ConvNeXt-Tiny, EfficientNet-B3, RegNetY-8GF, Swin-T, and ViT-B/16}---on four versions of the training data: the original images, images with randomly reassigned backgrounds, images with a fixed next-class background, and foreground-only images. These training conditions deliberately produce models with different degrees and directions of background dependence. This group therefore tests whether the same response semantics continue to track behavior as the model's learned use of background changes. Together, the two groups form a 72-model zoo: one provides models trained independently of our audit, while the other provides controlled diversity in the behavior being audited.

% ==================== 第一个大图（独占一行） ====================
\begin{figure*}[t]
\centering
\begin{subfigure}[t]{0.325\linewidth}
    \centering
    \includegraphics[width=\linewidth]{figures/imagenet9/fig5a_natural_behavioral_alignment.pdf}
    \caption{Behavioral alignment.}
\end{subfigure}\hfill
\begin{subfigure}[t]{0.34\linewidth}
    \centering
    \includegraphics[width=\linewidth]{figures/imagenet9/fig5b_matched_pair_mechanism_accuracy.pdf}
    \caption{Response-role accuracy.}
\end{subfigure}\hfill
\begin{subfigure}[t]{0.33\linewidth}
    \centering
    \includegraphics[width=\linewidth]{figures/imagenet9/fig5c_magnitude_stratified_macro.pdf}
    \caption{Across magnitude bins.}
\end{subfigure}

\caption{\textbf{Response semantics transfer to ImageNet-9.}
(a) Evidence and fragility align with independently measured natural-image behaviors.
(b) Nearly identical response magnitudes can correspond to different response semantics, which DECAF preserves.
(c) The advantage persists across magnitude bins.}
\label{fig:imagenet9-mechanism}
\end{figure*}

% ==================== 第二个大图（独占一行） ====================
\begin{figure*}[t]
\centering
\begin{subfigure}[t]{0.46\linewidth}
    \centering
    \includegraphics[width=\linewidth]{figures/imagenet9/fig6a_protocol_change_ratio.pdf}
    \caption{Response composition under patch reveal.}
\end{subfigure}
\hspace{5mm}
\begin{subfigure}[t]{0.39\linewidth}
    \centering
    \includegraphics[width=\linewidth]{figures/imagenet9/fig6b_protocol_rank_transfer.pdf}
    \caption{Model-rank transfer across reveal paths.}
\end{subfigure}

\caption{\textbf{Reveal paths change response composition and ranking.}
(a) Patch reveal increases ordinary magnitude mainly through contradiction and fragility rather than evidence.
(b) Ordinary-magnitude rankings transfer poorly across reveal paths, while evidence and contradiction remain substantially more stable.}
\label{fig:protocol-audit}
\end{figure*}

\subsection{Response semantics survive natural images}

The first question is whether the response roles validated in
Section~\ref{sec:controlled} remain connected to independently measured
behavior on natural images. We consider two behaviors. Background reliance
asks whether changing only the background substantially disrupts the model's
prediction. Endpoint-null sensitivity asks whether a pair with little final
background effect can nevertheless remain sensitive to held-out background
corruptions. Both behaviors are defined independently of \decaf. If the
response semantics transfer, $E$ should identify the former and $F$ the latter.

For background reliance, we mark a pair as positive when the model correctly classifies the image with a same-class background, but replacing that background with a random-class one either changes the predicted class or lowers the true-class probability by at least $0.20$. We then ask whether these behavior-positive pairs tend to receive larger response scores. Using $E$ directly as a ranking score gives an AUROC of $0.930$: a value of $0.5$ corresponds to random ranking and $1$ to perfect separation. Ordinary response magnitude reaches $0.900$, while endpoint magnitude reaches $0.960$ (Figure~\ref{fig:imagenet9-mechanism}(a)). The strong endpoint result is expected because this behavioral target is itself defined by a prediction change at the endpoint.

For endpoint-null sensitivity, we first restrict attention to pairs whose final background effect is negligible, $|d|<0.02$. We then apply a separate set of background corruptions that are not used to construct the reveal path. A pair is marked positive if any of these held-out corruptions changes the predicted class or shifts the true-class probability by at least $0.20$. We ask whether these independently identified sensitive pairs tend to receive larger $F$ values than the remaining endpoint-null pairs. $F$ reaches AUROC $0.878$, compared with $0.433$ for ordinary response magnitude, $0.316$ for endpoint magnitude, and $0.635$ for SmoothGrad (Figure~\ref{fig:imagenet9-mechanism}(a)).

Together, the two tests show complementary behavior: $E$ tracks consequential background use when the final effect is present, whereas $F$ exposes sensitivity that remains when that final effect is negligible. Appendix~\ref{app:imagenet9-setup} gives the full pair construction, held-out corruption set, behavioral-label definitions, and baseline details.


\subsection{Same Response Magnitude, Different Response Semantics}

We next ask whether the decomposition reveals information that ordinary
response magnitude does not already contain. We reuse the independently defined
evidence and fragility indicators from Section~6.1 and add a contradiction
indicator: contradiction behavior is present when changing only the background
makes the model switch from the foreground class to the class associated with
the new background. These three behavioral indicators may overlap.

We then compare cases with different behavioral patterns but nearly identical
ordinary response magnitudes. Requiring their $\Abs$ values to differ by at
most $5\%$ yields 8,289 matched comparisons. Because a case may exhibit more
than one behavior, we do not force it into a single ground-truth class.
Instead, we ask whether the largest of $E$, $C$, and $F$ corresponds to a
behavior that is actually present; overlap and ties are handled as detailed in
Appendix~\ref{app:imagenet9-setup}.

Under this test, a single magnitude provides no basis for preferring evidence,
contradiction, or fragility. $\Abs$ reaches a role-agreement accuracy of
$0.350$, whereas \decaf reaches $0.964$
(Figure~\ref{fig:imagenet9-mechanism}(b)). Thus, nearly identical response
magnitudes can accompany substantially different observed behaviors, while
their decomposition preserves this distinction.

The conclusion also holds beyond the matched cases. Within narrow bins of
ordinary response magnitude, macro-AUROC rises from $0.517$ for $\Abs$ to
$0.677$ for \decaf, and \decaf exceeds $\Abs$ in 15 of the 16 bins that
contain all three behavioral indicators
(Figure~\ref{fig:imagenet9-mechanism}(c)). Appendix~\ref{app:imagenet9-setup}
gives the exact matching rule, overlap and tie handling, per-bin support, and
complete baseline comparisons.

\subsection{What Changes When the Reveal Path Changes?}

We finally ask whether the same endpoint information can produce different
conclusions when it is revealed differently. We keep the model and both
factual--counterfactual endpoints fixed. In one path, the whole image gradually
changes from a common blurred image toward each endpoint; in the other, the
same endpoint information is revealed region by region. We call them the
blend and patch paths.

Ordinary response magnitude increases by about $1.8\times$ under the patch
path, but the increase is not uniform across response roles
(Figure~\ref{fig:protocol-audit}(a)). Evidence remains close to its blend value,
contradiction grows by about $1.8\times$, and fragility by more than
$4\times$. This is consistent with their semantics: with the endpoints fixed,
evidence remains oriented by the same final contrast, whereas fragility
explicitly measures sensitivity along the chosen reveal path. The pattern is
not guaranteed by the definition, but here the larger ordinary response comes
primarily from fragility rather than additional evidence.

We also ask whether the reveal path changes conclusions about how models
compare. For each model, we average $E$, $C$, $F$, and $\Abs$ over test images
and rank the models separately by each quantity under the blend and patch
paths. Spearman correlation measures the agreement between the two rankings.
Ordinary-magnitude rankings are highly unstable
($\rho=0.17/0.26$ for Same--Rand/Same--Next), whereas evidence rankings remain
at $0.86/0.77$ and contradiction rankings near $0.93$. Fragility is
intermediate at $0.69/0.71$, consistent with its greater path dependence
(Figure~\ref{fig:protocol-audit}(b)). Patch-order and threshold checks are
reported in Appendix~\ref{app:imagenet9-results}.